\begin{problem}[第30届全国中学生物理竞赛决赛][亚毫米细丝双光束干涉测径]{15}
亚毫米细丝直径的双光束干涉测量装置如图(a)所示，其中 $\mathrm{M}_{1} 、 \mathrm{M}_{2} 、 \mathrm{M}_{3}$ 为全反射镜，相机为 CCD (电荷耦合装置)相机。来自钒酸铱自倍频激光器的激光束被分束器分成两束，一束经全反镜 $\mathrm{M}_{1}$ 反射后从上侧入射到细丝上；另一束经全反镜 $\mathrm{M}_{3}$ 和 $\mathrm{M}_{2}$ 相继反射后从下侧入射到细丝上。图(b)给出了两条反射光线产生干涉的光路，其中 $\theta_{1}$ 、$ \theta_{2}$ 分别为上、下两侧的入射角， $D$ 为细丝轴线到观察屏(即相机感光片)的距离， $P$ 为两条反射线在观察屏上的交点。已建立这样的直角坐标系：坐标原点 $O$ 位于细丝的轴线上， $x$ 轴(未画出)沿细丝轴线、指向纸面内， $y$ 轴与入射到细丝上的光线平行 ( $y$ 轴的正向向上)， $z$ 轴指向观察屏并与其垂直。已知光波长为 $\lambda$ ，屏上干涉条纹的间距为 $a$ 。由于 $D$ 远大于细丝直径和观察屏的尺寸，可假设投射到屏上的只有非常接近平行于 $z$ 轴的细光束。
\begin{subproblem}
由于细丝到观察屏的距离远大于观察屏的尺寸，因而上、下两侧的入射光只有 $45^{\circ}$ 入射角附近的细光束经细丝反射到屏上，上、下两侧的反射光束分别形成两个虚像。试求这两个虚像的位置。(注：当 $x \sim 0$ 时， $\sin x \sim x, \cos x \sim 1$ )
\end{subproblem}
\begin{subproblem}
求细丝的直径 $d$ 。
\end{subproblem}
\end{problem}